%!TeX root=../main.tex

\chapter{Conclusioni}\label{ch:conclusioni}


\section{Sintesi del lavoro svolto e risultati ottenuti}\label{sec:sintesi-del-lavoro-svolto-e-risultati-ottenuti}
Il progetto di tesi ha portato alla progettazione, implementazione e distribuzione di \textbf{Swift Contest}, una piattaforma software completa e funzionante per la gestione di concorsi, che raggiunge con successo tutti gli obiettivi prefissati in fase di analisi.
È stato dimostrato come, attraverso l'integrazione di un ecosistema tecnologico moderno e accessibile, sia possibile costruire un sistema complesso, sicuro e scalabile, capace di risolvere le inefficienze e le criticità dei processi di gestione tradizionali.

I principali risultati ottenuti possono essere riassunti come segue:
\begin{itemize}
    \item \textbf{Realizzazione di un'applicazione cross-platform unificata}: utilizzando Flutter, è stato sviluppato un unico client che offre un'esperienza utente coerente e funzionale sia su piattaforme web che mobile (Android), validando l'efficacia di un approccio a codebase singola per ridurre i tempi di sviluppo e la complessità di manutenzione.
    \item \textbf{Implementazione di un'architettura backend sicura e robusta}: è stato costruito un backend basato sul Principio del Minimo Privilegio, dove l'accesso ai dati è interamente mediato da un'API controllata, proteggendo l'integrità del sistema.
    \item \textbf{Creazione di un sistema di votazione affidabile e trasparente}: il meccanismo di \emph{snapshot} delle sessioni di voto garantisce l'immutabilità dei dati e l'auditabilità del processo, risolvendo una delle principali criticità dei sistemi di valutazione manuali.
    \item \textbf{Validazione di un'infrastruttura efficiente e sostenibile}: l'architettura basata su servizi cloud (Supabase, Vercel) e l'uso mirato di risorse onerose come il \emph{realtime} dimostrano la fattibilità di costruire soluzioni complesse mantenendo i costi operativi contenuti, in linea con i piani gratuiti o a basso costo offerti dai provider.
\end{itemize}
Il prototipo realizzato non è solo una dimostrazione tecnica, ma si presenta come un prodotto maturo e coeso, già utilizzabile in contesti reali a piccola e media scala.

\section{Limiti e analisi critica}\label{sec:limiti-e-analisi-critica}
Un'analisi critica e onesta del lavoro svolto è fondamentale per comprenderne appieno il valore e i margini di miglioramento.
Nonostante i risultati positivi, il progetto presenta alcuni limiti, intrinseci a un elaborato sviluppato in un contesto accademico con vincoli temporali definiti.
\begin{itemize}
    \item \textbf{Logica di ranking semplificata}: l'attuale meccanismo di calcolo delle classifiche, basato sulla semplice somma dei punteggi dei campi \emph{slider}, è funzionale ma basilare.
    Manca la possibilità di introdurre regole più complesse, come l'assegnazione di pesi diversi ai campi di voto o l'applicazione di algoritmi di scarto per i voti estremi (es.
    eliminare il voto più alto e più basso), limitandone l'applicabilità in contesti con regolamenti di valutazione più sofisticati.
    \item \textbf{Assenza di versioning dei form di voto}: la stabilità delle votazioni è garantita tramite snapshot, ma non esiste un vero sistema di versionamento per i \emph{VotingForm}.
    Una volta che un form è stato usato, modificarlo potrebbe creare incongruenze se lo si volesse riutilizzare in futuro.
    Una gestione delle versioni permetterebbe agli organizzatori di creare e archiviare template di form riutilizzabili.
    \item \textbf{Osservabilità e monitoraggio limitati}: il sistema non integra strumenti avanzati di logging, metriche o tracing distribuito.
    In un ambiente di produzione su larga scala, l'assenza di questi strumenti renderebbe più complessa l'analisi delle performance, il debug di errori specifici e il monitoraggio proattivo della salute del sistema.
    \item \textbf{Test automatici non esaustivi}: sebbene la logica di business nei BLoC sia stata coperta da test di unità, il progetto beneficerebbe enormemente dall'aggiunta di una suite di test più completa, includendo test di integrazione per verificare l'interazione tra client e backend e test end-to-end per simulare i flussi utente completi.
\end{itemize}

\section{Prospettive e sviluppi futuri}\label{sec:prospettive-e-sviluppi-futuri}
Il progetto attuale costituisce una solida base per numerosi e interessanti sviluppi futuri, che potrebbero trasformare Swift Contest da un prototipo avanzato a un prodotto commerciale completo e competitivo.
\begin{itemize}
    \item \textbf{Ranking avanzati e configurabili}: lo sviluppo più immediato e di maggior valore sarebbe l'introduzione di un sistema che permetta agli organizzatori di definire regole di calcolo personalizzate per le classifiche, aumentando la flessibilità della piattaforma.
    \item \textbf{Supporto per sottomissioni multiple (Multi-Work)}: estendere il modello dati per permettere ai partecipanti di inviare più di un \emph{Work} all'interno dello stesso concorso è una funzionalità molto richiesta, specialmente in ambiti come la fotografia o la poesia.
    \item \textbf{Integrazione di notifiche push}: affiancare alle notifiche inbox un sistema di notifiche push (es.\ tramite Firebase Cloud Messaging o un altro provider) aumenterebbe significativamente il coinvolgimento e la reattività degli utenti, notificandoli istantaneamente di eventi importanti.
    \item \textbf{Hardening della sicurezza}: sebbene l'architettura sia sicura, potrebbe essere ulteriormente rafforzata con meccanismi di \emph{rate limiting} sulle Edge Functions per prevenire abusi e attacchi di tipo DoS, e con policy RLS ancora più granulari.
    \item \textbf{Internazionalizzazione (i18n) e accessibilità (a11y)}: la traduzione dell'interfaccia utente in più lingue e il miglioramento del supporto per le tecnologie assistive (es.\ screen reader) sono passi fondamentali per rendere l'applicazione accessibile a un pubblico globale e inclusivo.
    \item \textbf{Distribuzione ufficiale su App Store}: per completare la visione cross-platform, un passo futuro sarebbe quello di creare una versione per iOS e pubblicare entrambe le applicazioni mobile sugli store ufficiali (Google Play Store e Apple App Store), semplificandone la distribuzione e gli aggiornamenti.
\end{itemize}

\section{Considerazioni finali}\label{sec:considerazioni-finali}
Swift Contest rappresenta un esempio concreto di come un'architettura software ben ponderata, unita a un ecosistema tecnologico moderno e accessibile, possa portare alla creazione di un sistema complesso, sicuro e di valore.
La scelta di combinare la flessibilità di Flutter con la potenza e la sicurezza di Supabase si è rivelata vincente, permettendo di raggiungere un ottimo equilibrio tra rapidità di sviluppo, ricchezza funzionale e sostenibilità dei costi.

Il progetto non si esaurisce con la presente tesi; esso costituisce una base di partenza robusta e scalabile per future iterazioni e, potenzialmente, per applicazioni in contesti reali.
Il lavoro svolto conferma la validità dell'approccio architetturale adottato e apre prospettive interessanti nel campo delle piattaforme collaborative e dei sistemi di votazione digitale, dimostrando che è possibile realizzare prototipi di alta qualità anche con le risorse a disposizione di uno studente o di un piccolo team di sviluppo.